\documentclass[DIN, pagenumber=false, fontsize=11pt, parskip=half]{scrartcl}
\usepackage[UTF8]{ctex}  % 中文支持
% \usepackage{ngerman}
\usepackage[utf8]{inputenc}
\usepackage[T1]{fontenc}
\usepackage{textcomp}
\usepackage{xcolor} % Required for custom red
\usepackage{amsmath, amssymb} % 数学公式必备
\usepackage{geometry} % Custom margins
\usepackage{enumitem}
\usepackage{tocloft} % Modifying the table of contents

% for matlab code
% bw = blackwhite - optimized for print, otherwise source is colored
\usepackage[framed,numbered,bw]{mcode}
\usepackage[most]{tcolorbox}
% for other code
%\usepackage{listings}

\setlength{\parindent}{0em}

% set section in CM
\setkomafont{section}{\normalfont\bfseries\Large}

\newcommand{\mytitle}[1]{{\noindent\Large\textbf{#1}}}
\newcommand{\mysection}[1]{\textbf{\section*{#1}}}
\newcommand{\mysubsection}[2]{\romannumeral #1) #2}

\newtcolorbox{parchmentquote}{
    enhanced,
    colback=orange!5!white,   % 浅米黄色背景
    colframe=orange!50!black,
    boxrule=0.5pt,
    arc=2pt,                  % 细微圆角
    boxsep=5pt,
    before skip=10pt,
    after skip=10pt,
    fontupper=\itshape,
    borderline west={2pt}{0pt}{orange!80!black} % 加厚的左装饰线
}
% 设置选项样式：使用 A, B, C, D

%===================================
\begin{document}

\noindent\textbf{课下练习} \hfill \textbf{online course}\\
primary school \hfill Liu\\

\mytitle{分数加法减法 \hfill \today}
\newline
\begin{parchmentquote}
	难点在异分母分数加减法
\end{parchmentquote}


%===================================
\section{选择题}
下面4组数中，()组中的数都是质数。
\begin{enumerate}[label=\Alph*.]
  \item $13,71,51$
  \item $47,79,93$
  \item $31,73,97$
  \item $2,17,91$
\end{enumerate}
当a是自然数时，4a+1一定是()。

\begin{enumerate}[label=\Alph*.]
  \item 奇数 \item 偶数
  \item 质数
  \item 偶数
\end{enumerate}
用2、3、5、8四张卡片摆数，摆出的四位数()。

\begin{enumerate}[label=\Alph*.]
  \item 一定是2的倍数
  \item 一定是5的倍数
  \item 一定是3的倍数
  \item 不确定是不是2、3、5的倍数
\end{enumerate}
如果$\frac{12}{18}$的分子减去4，要使分数的大小不变，分母可以()。

\begin{enumerate}[label=\Alph*.]
  \item 减去4
  \item 减去15
  \item 除以4
  \item 除以5
\end{enumerate}

\section{计算题}
\begin{enumerate}
    \item \(\frac{2}{11} + \frac{3}{11} = \underline{\qquad}\)
    \item \(\frac{11}{12} + \frac{13}{24} = \underline{\qquad}\)
    \item \(\frac{2}{3} - \frac{5}{6}+\frac{1}{3} = \underline{\qquad}\)
    \item \(0.8 + \frac{3}{5} = \underline{\qquad}\)
    \item \(2.25 - \frac{3}{4}+ \frac{2}{3}  = \underline{\qquad}\)
    \item \(17 - 1\frac{17}{19}- \frac{2}{19}  = \underline{\qquad}\)
    \item \(\frac{5}{6}- \frac{2}{9}  = \underline{\qquad}\)
   
\end{enumerate}



\end{document}